\documentclass{article}
\usepackage[utf8]{inputenc}
\usepackage[vietnamese]{babel}
\usepackage{amsmath}
\usepackage{amssymb}
\usepackage{enumitem}
\usepackage{geometry}
\geometry{a4paper, margin=1in}

\title{Ngân hàng câu hỏi Toán Lớp 4 (Crawl từ Vietjack)}
\author{Hệ thống Ôn luyện Toán Bổ Trợ}
\date{\today}

\begin{document}
\maketitle

\tableofcontents
\newpage
\section{Lớp 4}
\subsection{Chủ đề 1. Số tự nhiên: 6. Bài 6: Các số có nhiều chữ số}
\begin{minipage}{\textwidth}
\noindent\textbf{Câu hỏi 1 (MEDIUM):} Số 123 678 được đọc như thế nào? 
\begin{itemize}[label={}]
    \item A. Một trăm hai ba nghìn sáu trăm bảy tám.
    \item B. Một trăm hai mươi ba nghìn sáu trăm bảy tám.
    \item C. Một trăm hai mươi ba nghìn sáu trăm bảy mươi tám.
    \item D. Một trăm hai ba nghìn sáu trăm bảy mươi tám.
\end{itemize}
\textbf{Đáp án đúng:} C \\
\textbf{Lời giải chi tiết:} Số 123 678 được đọc là: Một trăm hai mươi ba nghìn sáu trăm bảy mươi tám. \\
\textbf{Lỗi thường gặp:}
\begin{itemize}
  \item \textbf{Đáp án A (Nhầm lẫn số mươi và nghìn):} Học sinh nhầm lẫn số mươi và nghìn, đọc 'hai ba' là 'hai mươi ba' và bỏ qua số nghìn.
  \item \textbf{Đáp án B (Nhầm lẫn số mươi và nghìn):} Học sinh nhầm lẫn số mươi và nghìn, đọc 'hai ba' là 'hai mươi ba' và bỏ qua số nghìn.
  \item \textbf{Đáp án D (Nhầm lẫn số mươi và nghìn):} Học sinh nhầm lẫn số mươi và nghìn, đọc 'hai ba' là 'hai mươi ba' và bỏ qua số nghìn.
\end{itemize}
\end{minipage}
\vspace{1em}

\begin{minipage}{\textwidth}
\noindent\textbf{Câu hỏi 2 (MEDIUM):} Số “Bảy mươi triệu” có bao nhiêu chữ số? 
\begin{itemize}[label={}]
    \item A. 4
    \item B. 6
    \item C. 8
    \item D. 9
\end{itemize}
\textbf{Đáp án đúng:} C \\
\textbf{Lời giải chi tiết:} Số “Bảy mươi triệu” được viết là: 70 000 000. Vậy số “Bảy mươi triệu” có: 8 chữ số. \\
\textbf{Lỗi thường gặp:}
\begin{itemize}
  \item \textbf{Đáp án A (Nhầm lẫn số hàng):} Học sinh nhầm lẫn số hàng triệu với số hàng nghìn, dẫn đến việc đếm sai số chữ số.
  \item \textbf{Đáp án B (Nhầm lẫn số hàng):} Học sinh nhầm lẫn số hàng triệu với số hàng nghìn, dẫn đến việc đếm sai số chữ số.
  \item \textbf{Đáp án D (Nhầm lẫn số hàng):} Học sinh nhầm lẫn số hàng triệu với số hàng nghìn, dẫn đến việc đếm sai số chữ số.
\end{itemize}
\end{minipage}
\vspace{1em}

\begin{minipage}{\textwidth}
\noindent\textbf{Câu hỏi 3 (MEDIUM):} Số “ Một trăm triệu” có bao nhiêu chữ số 0? 
\begin{itemize}[label={}]
    \item A. 5
    \item B. 7
    \item C. 8
    \item D. 9
\end{itemize}
\textbf{Đáp án đúng:} C \\
\textbf{Lời giải chi tiết:} Số “ Một trăm triệu” được viết là: 100 000 000. Vậy số “ Một trăm triệu” có 8 chữ số 0. \\
\textbf{Lỗi thường gặp:}
\begin{itemize}
  \item \textbf{Đáp án A (Nhầm lẫn vị trí chữ số 0):} Học sinh có thể nhầm lẫn vị trí chữ số 0 với số 1, dẫn đến việc đếm sai số 0. Ví dụ, họ có thể nghĩ rằng số 100 000 000 có 5 chữ số 0 vì họ đã đếm từ số 1 đến số 5, bỏ qua số 0 ở giữa.
  \item \textbf{Đáp án B (Tính sai số 0):} Học sinh có thể tính sai số 0 vì họ đã đếm sai số 0. Ví dụ, họ có thể nghĩ rằng số 100 000 000 có 7 chữ số 0 vì họ đã đếm từ số 1 đến số 7, bỏ qua số 0 ở giữa.
  \item \textbf{Đáp án D (Nhầm lẫn số 0 với số khác):} Học sinh có thể nhầm lẫn số 0 với số khác, dẫn đến việc đếm sai số 0. Ví dụ, họ có thể nghĩ rằng số 100 000 000 có 9 chữ số 0 vì họ đã nhầm lẫn số 0 với số 9.
\end{itemize}
\end{minipage}
\vspace{1em}

\begin{minipage}{\textwidth}
\noindent\textbf{Câu hỏi 4 (MEDIUM):} Số “Hai mươi ba triệu” viết là: 
\begin{itemize}[label={}]
    \item A. 23 000
    \item B. 230 000
    \item C. 2 300 000
    \item D. 23 000 000
\end{itemize}
\textbf{Đáp án đúng:} D \\
\textbf{Lời giải chi tiết:} Số “Hai mươi ba triệu” viết là: 23 000 000. \\
\textbf{Lỗi thường gặp:}
\begin{itemize}
  \item \textbf{Đáp án A (Lỗi tính toán):} Học sinh chọn đáp án A do nhầm lẫn trong các bước cộng, trừ, nhân, chia hoặc đọc sai đơn vị đo.
  \item \textbf{Đáp án B (Lỗi tính toán):} Học sinh chọn đáp án B do nhầm lẫn trong các bước cộng, trừ, nhân, chia hoặc đọc sai đơn vị đo.
  \item \textbf{Đáp án C (Lỗi tính toán):} Học sinh chọn đáp án C do nhầm lẫn trong các bước cộng, trừ, nhân, chia hoặc đọc sai đơn vị đo.
\end{itemize}
\end{minipage}
\vspace{1em}

\begin{minipage}{\textwidth}
\noindent\textbf{Câu hỏi 5 (MEDIUM):} Phát biểu sau đây đúng hay sai? \\
“Số nhỏ nhất có 8 chữ số là 90 000 000” 
\begin{itemize}[label={}]
    \item A. Đúng.
    \item B. Sai.
\end{itemize}
\textbf{Đáp án đúng:} B \\
\textbf{Lời giải chi tiết:} Số nhỏ nhất có 8 chữ số là: 10 000 000. Vậy phát biểu trên là: Sai. \\
\textbf{Lỗi thường gặp:}
\begin{itemize}
  \item \textbf{Đáp án A (Nhầm lẫn số 0):} Học sinh có thể nhầm lẫn số 0 với số 9, dẫn đến việc tính ra số 90 000 000 thay vì số 10 000 000.
  \item \textbf{Đáp án B (Nhầm lẫn vị trí số 9):} Học sinh có thể nhầm lẫn vị trí số 9 với số 0, dẫn đến việc tính ra số 10 000 000 thay vì số 1 000 000 000.
\end{itemize}
\end{minipage}
\vspace{1em}

\begin{minipage}{\textwidth}
\noindent\textbf{Câu hỏi 6 (MEDIUM):} Số cần điền là: \\
1 200 000; .....; 1 400 000; 1 500 000 
\begin{itemize}[label={}]
    \item A. 1 300 000.
    \item B. 1 250 000.
    \item C. 1 350 000.
    \item D. 1 520 000.
\end{itemize}
\textbf{Đáp án đúng:} A \\
\textbf{Lời giải chi tiết:} Dãy số trên gồm các số cách nhau 100 000 đơn vị. Vậy số cần điền là: 1 200 000 + 100 000 = 1 300 000. \\
\textbf{Lỗi thường gặp:}
\begin{itemize}
  \item \textbf{Đáp án B (Nhầm lẫn số lần nhân):} Học sinh có thể nhầm lẫn số lần nhân 100 000, dẫn đến việc chọn đáp án 1 250 000. Họ có thể tính như sau: 1 200 000 + 50 000 = 1 250 000, nhưng thực tế số cần điền là 1 300 000.
  \item \textbf{Đáp án C (Nhầm đơn vị đo):} Học sinh có thể nhầm lẫn đơn vị đo, dẫn đến việc chọn đáp án 1 350 000. Họ có thể tính như sau: 1 200 000 + 150 000 = 1 350 000, nhưng thực tế số cần điền là 1 300 000.
  \item \textbf{Đáp án D (Nhầm lẫn số lần nhân và đơn vị đo):} Học sinh có thể nhầm lẫn số lần nhân và đơn vị đo, dẫn đến việc chọn đáp án 1 520 000. Họ có thể tính như sau: 1 500 000 + 20 000 = 1 520 000, nhưng thực tế số cần điền là 1 300 000.
\end{itemize}
\end{minipage}
\vspace{1em}

\begin{minipage}{\textwidth}
\noindent\textbf{Câu hỏi 7 (MEDIUM):} Số có năm chữ số 0 và số có năm chữ số, số nào lớn hơn? 
\begin{itemize}[label={}]
    \item A. Số có năm chữ số 0.
    \item B. Số có năm chữ số.
    \item C. Hai số bằng nhau.
    \item D. Không xác định được.
\end{itemize}
\textbf{Đáp án đúng:} A \\
\textbf{Lời giải chi tiết:} Số có năm chữ số 0 phải là số ít nhất có sáu chữ số. Vậy số có năm chữ số 0 lớn hơn số có năm chữ số. \\
\textbf{Lỗi thường gặp:}
\begin{itemize}
  \item \textbf{Đáp án B (Nhầm lẫn về số lượng chữ số):} Học sinh có thể nhầm lẫn giữa số có năm chữ số và số có năm chữ số 0. Họ có thể nghĩ rằng số có năm chữ số lớn hơn số có năm chữ số 0 vì số có năm chữ số có thể có nhiều chữ số hơn. Tuy nhiên, thực tế là số có năm chữ số 0 phải là số ít nhất có sáu chữ số, vì vậy nó lớn hơn số có năm chữ số.
  \item \textbf{Đáp án C (Nhầm lẫn về so sánh):} Học sinh có thể nghĩ rằng hai số bằng nhau vì chúng đều có năm chữ số. Tuy nhiên, thực tế là số có năm chữ số 0 phải là số ít nhất có sáu chữ số, vì vậy nó lớn hơn số có năm chữ số.
  \item \textbf{Đáp án D (Nhầm lẫn về không xác định):} Học sinh có thể nghĩ rằng không thể xác định được số lớn hơn vì chúng đều có năm chữ số. Tuy nhiên, thực tế là số có năm chữ số 0 phải là số ít nhất có sáu chữ số, vì vậy nó lớn hơn số có năm chữ số.
\end{itemize}
\end{minipage}
\vspace{1em}

\begin{minipage}{\textwidth}
\noindent\textbf{Câu hỏi 8 (MEDIUM):} Mười tờ tiền có mệnh giá hai mươi nghìn thì tổng số tiền là: 
\begin{itemize}[label={}]
    \item A. 20 000 đồng.
    \item B. 200 000 đồng.
    \item C. 2 000 đồng.
    \item D. 2 000 000 đồng.
\end{itemize}
\textbf{Đáp án đúng:} B \\
\textbf{Lời giải chi tiết:} Tổng số tiền là: 20 000 × 10 = 200 000 (đồng) Đáp số: 200 000 đồng \\
\textbf{Lỗi thường gặp:}
\begin{itemize}
  \item \textbf{Đáp án A (Nhầm lẫn đơn vị đo):} Học sinh có thể nhầm lẫn đơn vị đo giữa nghìn và triệu, dẫn đến việc tính ra đáp án 20 000 đồng. Để tránh nhầm lẫn này, học sinh cần đọc kỹ đề bài và đảm bảo rằng họ đang sử dụng đơn vị đo đúng.
  \item \textbf{Đáp án C (Tính sai số):} Học sinh có thể tính sai số khi nhân 20 000 với 10, dẫn đến việc tính ra đáp án 2 000 đồng. Để tránh sai số này, học sinh cần kiểm tra lại công thức và đảm bảo rằng họ đang sử dụng số liệu chính xác.
  \item \textbf{Đáp án D (Nhầm lẫn số mũ):} Học sinh có thể nhầm lẫn số mũ giữa 10 và 1 triệu, dẫn đến việc tính ra đáp án 2 000 000 đồng. Để tránh nhầm lẫn này, học sinh cần đọc kỹ đề bài và đảm bảo rằng họ đang sử dụng số mũ đúng.
\end{itemize}
\end{minipage}
\vspace{1em}

\begin{minipage}{\textwidth}
\noindent\textbf{Câu hỏi 9 (MEDIUM):} Số lẻ nhỏ nhất có bảy chữ số là? 
\begin{itemize}[label={}]
    \item A. 1 000 010.
    \item B. 1 000 001.
    \item C. 1 000 009.
    \item D. 9 000 001
\end{itemize}
\textbf{Đáp án đúng:} B \\
\textbf{Lời giải chi tiết:} Số lẻ nhỏ nhất có bảy chữ số là: 1 000 001. \\
\textbf{Lỗi thường gặp:}
\begin{itemize}
  \item \textbf{Đáp án A (Nhầm lẫn vị trí chữ số 0):} Học sinh có thể nhầm lẫn vị trí chữ số 0 trong số 1 000 010, dẫn đến việc chọn đáp án A. Để tránh nhầm lẫn này, học sinh cần đọc kỹ và hiểu rõ vị trí của các chữ số trong số lẻ nhỏ nhất có bảy chữ số.
  \item \textbf{Đáp án C (Nhầm lẫn với số chẵn):} Học sinh có thể nhầm lẫn số lẻ 1 000 001 với số chẵn 1 000 000, dẫn đến việc chọn đáp án C. Để tránh nhầm lẫn này, học sinh cần hiểu rõ đặc điểm của số lẻ và số chẵn.
  \item \textbf{Đáp án D (Nhầm lẫn vị trí chữ số 9):} Học sinh có thể nhầm lẫn vị trí chữ số 9 trong số 9 000 001, dẫn đến việc chọn đáp án D. Để tránh nhầm lẫn này, học sinh cần đọc kỹ và hiểu rõ vị trí của các chữ số trong số lẻ nhỏ nhất có bảy chữ số.
\end{itemize}
\end{minipage}
\vspace{1em}

\begin{minipage}{\textwidth}
\noindent\textbf{Câu hỏi 10 (MEDIUM):} Mười chục triệu còn được gọi là: 
\begin{itemize}[label={}]
    \item A. Một trăm triệu.
    \item B. Một triệu.
    \item C. Một triệu triệu.
    \item D. Mười triệu trăm.
\end{itemize}
\textbf{Đáp án đúng:} A \\
\textbf{Lời giải chi tiết:} Mười chục triệu còn được gọi là một trăm triệu. \\
\textbf{Lỗi thường gặp:}
\begin{itemize}
  \item \textbf{Đáp án B (Nhầm lẫn đơn vị đo):} Học sinh có thể nhầm lẫn giữa đơn vị triệu và đơn vị triệu triệu, dẫn đến việc chọn đáp án B. Một triệu. Để tránh nhầm lẫn này, học sinh cần đọc kỹ và hiểu rõ ý nghĩa của mỗi đơn vị đo.
  \item \textbf{Đáp án C (Tính toán sai):} Học sinh có thể tính toán sai khi đọc từ 'mười triệu' thành 'một triệu triệu'. Để tránh nhầm lẫn này, học sinh cần đọc từ và tính toán cẩn thận.
  \item \textbf{Đáp án D (Nhầm lẫn cấu trúc từ):} Học sinh có thể nhầm lẫn cấu trúc từ 'mười triệu' thành 'mười triệu trăm'. Để tránh nhầm lẫn này, học sinh cần đọc từ và hiểu rõ ý nghĩa của mỗi từ.
\end{itemize}
\end{minipage}
\vspace{1em}

\subsection{Chủ đề 1. Số tự nhiên: 7. Bài 7: Các số có nhiều chữ số (tiếp theo)}
\begin{minipage}{\textwidth}
\noindent\textbf{Câu hỏi 1 (MEDIUM):} Lớp nghìn có hàng nào? 
\begin{itemize}[label={}]
    \item A. Trăm nghìn.
    \item B. Chục nghìn.
    \item C. Nghìn.
    \item D. Trăm nghìn, chục nghìn, nghìn.
\end{itemize}
\textbf{Đáp án đúng:} D \\
\textbf{Lời giải chi tiết:} Lớp nghìn có hàngtrăm nghìn, chục nghìn, nghìn. \\
\textbf{Lỗi thường gặp:}
\begin{itemize}
  \item \textbf{Đáp án A (Nhầm lẫn hàng nghìn):} Học sinh nhầm lẫn giữa hàng nghìn và hàng trăm nghìn, nghĩ rằng hàng nghìn chỉ bao gồm hàng trăm nghìn.
  \item \textbf{Đáp án B (Nhầm lẫn đơn vị):} Học sinh nhầm lẫn giữa đơn vị nghìn và đơn vị chục nghìn, nghĩ rằng hàng nghìn chỉ bao gồm hàng chục nghìn.
\end{itemize}
\end{minipage}
\vspace{1em}

\begin{minipage}{\textwidth}
\noindent\textbf{Câu hỏi 2 (MEDIUM):} Chữ số 2 trong số 137 287 341 thuộc hàng nào? 
\begin{itemize}[label={}]
    \item A. Nghìn.
    \item B. Chục nghìn.
    \item C. Trăm nghìn.
    \item D. Triệu.
\end{itemize}
\textbf{Đáp án đúng:} C \\
\textbf{Lời giải chi tiết:} Chữ số 2 trong số 137 287 341 thuộc hàngtrăm nghìn, lớp nghìn. \\
\textbf{Lỗi thường gặp:}
\begin{itemize}
  \item \textbf{Đáp án A (Nhầm đơn vị đo):} Học sinh có thể nhầm lẫn chữ số 2 thuộc hàng nghìn, vì số 3 ở hàng nghìn có thể khiến họ nghĩ rằng chữ số 2 cũng thuộc hàng nghìn. Để tránh nhầm lẫn này, học sinh cần đọc số từ trái sang phải và xác định hàng của mỗi chữ số.
  \item \textbf{Đáp án B (Nhầm lẫn vị trí chữ số):} Học sinh có thể nhầm lẫn vị trí của chữ số 2, nghĩ rằng nó thuộc hàng chục nghìn vì số 3 ở hàng chục nghìn. Để tránh nhầm lẫn này, học sinh cần đọc số từ trái sang phải và xác định hàng của mỗi chữ số.
  \item \textbf{Đáp án D (Nhầm đơn vị đo):} Học sinh có thể nhầm lẫn chữ số 2 thuộc hàng triệu, vì số 1 ở hàng triệu có thể khiến họ nghĩ rằng chữ số 2 cũng thuộc hàng triệu. Để tránh nhầm lẫn này, học sinh cần đọc số từ trái sang phải và xác định hàng của mỗi chữ số.
\end{itemize}
\end{minipage}
\vspace{1em}

\begin{minipage}{\textwidth}
\noindent\textbf{Câu hỏi 3 (MEDIUM):} Lớp đơn vị của số 67 214 236 gồm những số nào? 
\begin{itemize}[label={}]
    \item A. 6 và 7
    \item B. 2, 1, 4
    \item C. 2, 3, 6
    \item D. 7, 4, 6
\end{itemize}
\textbf{Đáp án đúng:} C \\
\textbf{Lời giải chi tiết:} Số 67 214 236 có: Lớp đơn vị: 2, 3, 6. \\
\textbf{Lỗi thường gặp:}
\begin{itemize}
  \item \textbf{Đáp án A (Nhầm lẫn chữ số với lớp đơn vị):} Học sinh nhầm lẫn chữ số 6 và 7 với lớp đơn vị, dẫn đến việc chọn đáp án A. Để làm đúng, học sinh cần tách rời chữ số và lớp đơn vị để xác định lớp đơn vị chính xác.
  \item \textbf{Đáp án B (Nhầm lẫn vị trí chữ số):} Học sinh nhầm lẫn vị trí của chữ số 2, 1, 4 và chọn đáp án B. Để làm đúng, học sinh cần xác định vị trí của mỗi chữ số để tính lớp đơn vị.
  \item \textbf{Đáp án D (Nhầm lẫn chữ số với lớp đơn vị):} Học sinh nhầm lẫn chữ số 7, 4, 6 với lớp đơn vị, dẫn đến việc chọn đáp án D. Để làm đúng, học sinh cần tách rời chữ số và lớp đơn vị để xác định lớp đơn vị chính xác.
\end{itemize}
\end{minipage}
\vspace{1em}

\begin{minipage}{\textwidth}
\noindent\textbf{Câu hỏi 4 (MEDIUM):} Giá trị của chữ số 7 trong số 700 312 là bao nhiêu? 
\begin{itemize}[label={}]
    \item A. 700 000
    \item B. 70 000
    \item C. 7 000
    \item D. 700
\end{itemize}
\textbf{Đáp án đúng:} A \\
\textbf{Lời giải chi tiết:} Giá trị của chữ số 7 trong số 700 312 là:700 000. \\
\textbf{Lỗi thường gặp:}
\begin{itemize}
  \item \textbf{Đáp án B (Nhầm lẫn vị trí chữ số):} Học sinh nhầm lẫn vị trí chữ số 7 với chữ số 0, dẫn đến việc tính ra đáp án 70 000.
  \item \textbf{Đáp án C (Nhầm lẫn độ lớn):} Học sinh không hiểu rõ về độ lớn của chữ số 7 trong số 700 312, dẫn đến việc tính ra đáp án 7 000.
  \item \textbf{Đáp án D (Nhầm lẫn đơn vị):} Học sinh nhầm lẫn đơn vị của chữ số 7, dẫn đến việc tính ra đáp án 700.
\end{itemize}
\end{minipage}
\vspace{1em}

\begin{minipage}{\textwidth}
\noindent\textbf{Câu hỏi 5 (MEDIUM):} Trong các số sau, số nào có chữ số hàng triệu là 7? 
\begin{itemize}[label={}]
    \item A. 34 700 000
    \item B. 27 435 000
    \item C. 4 079 800
    \item D. 126 798 001
    \item A. Số 34 700 000 có chữ số 7 thuộc hàng trăm nghìn.
    \item B. Số 27 435 000 có chữ số 7 thuộc hàng triệu.
    \item C. Số 4 079 800 có chữ số 7 thuộc hàng chục nghìn.
    \item D. Số 126 798 001 có chữ số 7 thuộc hàng trăm nghìn.
\end{itemize}
\textbf{Đáp án đúng:} B \\
\textbf{Lời giải chi tiết:} A.Số 34 700 000 có chữ số 7 thuộc hàng trăm nghìn. B.Số 27 435 000 có chữ số 7 thuộc hàng triệu. C.Số 4 079 800 có chữ số 7 thuộc hàng chục nghìn. D.Số 126 798 001 có chữ số 7 thuộc hàng trăm nghìn. \\
\textbf{Lỗi thường gặp:}
\begin{itemize}
  \item \textbf{Đáp án A (Nhầm lẫn vị trí chữ số 7):} Học sinh nhầm lẫn vị trí chữ số 7 trong số 34 700 000, nghĩ rằng nó thuộc hàng triệu khi thực tế nó thuộc hàng trăm nghìn.
  \item \textbf{Đáp án C (Nhầm lẫn vị trí chữ số 7):} Học sinh nhầm lẫn vị trí chữ số 7 trong số 4 079 800, nghĩ rằng nó thuộc hàng triệu khi thực tế nó thuộc hàng chục nghìn.
  \item \textbf{Đáp án D (Nhầm lẫn vị trí chữ số 7):} Học sinh nhầm lẫn vị trí chữ số 7 trong số 126 798 001, nghĩ rằng nó thuộc hàng triệu khi thực tế nó thuộc hàng trăm nghìn.
\end{itemize}
\end{minipage}
\vspace{1em}

\begin{minipage}{\textwidth}
\noindent\textbf{Câu hỏi 6 (MEDIUM):} Số 18 342 964 có chữ số 3 thuộc hàng .... lớp ...? 
\begin{itemize}[label={}]
    \item A. Nghìn, nghìn.
    \item B. Chục nghìn, nghìn.
    \item C. Trăm nghìn, nghìn.
    \item D. Triệu, triệu.
\end{itemize}
\textbf{Đáp án đúng:} C \\
\textbf{Lời giải chi tiết:} Số 18 342 964 có chữ số 3 thuộc hàngtrăm nghìn, lớp nghìn. \\
\textbf{Lỗi thường gặp:}
\begin{itemize}
  \item \textbf{Đáp án A (Nhầm lẫn đơn vị đo):} Học sinh nhầm lẫn giữa đơn vị 'nghìn' và 'trăm nghìn', dẫn đến việc chọn đáp án A. Để làm đúng, học sinh cần chú ý đến vị trí của chữ số 3 trong số 18 342 964, đó là hàng trăm nghìn.
  \item \textbf{Đáp án B (Nhầm lẫn vị trí chữ số):} Học sinh nhầm lẫn vị trí của chữ số 3, nghĩ rằng nó thuộc hàng chục nghìn, nhưng thực tế nó thuộc hàng trăm nghìn. Để làm đúng, học sinh cần đọc số từ trái sang phải và xác định vị trí của chữ số 3.
  \item \textbf{Đáp án D (Nhầm lẫn đơn vị đo và vị trí chữ số):} Học sinh nhầm lẫn giữa đơn vị 'triệu' và vị trí của chữ số 3, nghĩ rằng nó thuộc hàng triệu, nhưng thực tế nó thuộc hàng trăm nghìn. Để làm đúng, học sinh cần chú ý đến vị trí của chữ số 3 và đơn vị đo phù hợp.
\end{itemize}
\end{minipage}
\vspace{1em}

\begin{minipage}{\textwidth}
\noindent\textbf{Câu hỏi 7 (MEDIUM):} Số 6 124 840 được viết thành tổng là 
\begin{itemize}[label={}]
    \item A. 6 000 000 + 100 000 + 20 000 + 4 000 + 800 + 40 + 10.
    \item B. 6 000 000 + 100 000 + 20 000 + 4 000 + 800 + 40.
    \item C. 6 000 000 + 100 000 + 20 000 + 4 000 + 800 + 4.
    \item D. 6 000 000 + 100 000 + 20 000 + 4 000 + 80 + 40.
\end{itemize}
\textbf{Đáp án đúng:} B \\
\textbf{Lời giải chi tiết:} Ta có: 6 124 840 =6 000 000 + 100 000 + 20 000 + 4 000 + 800 + 40. \\
\textbf{Lỗi thường gặp:}
\begin{itemize}
  \item \textbf{Đáp án A (Nhầm lẫn số thứ tự):} Học sinh đã nhầm lẫn số thứ tự của các chữ số trong số 6 124 840, dẫn đến việc thêm 10 vào cuối tổng.
  \item \textbf{Đáp án C (Nhầm lẫn đơn vị đo):} Học sinh đã nhầm lẫn đơn vị đo của số 800, dẫn đến việc thêm 4 vào cuối tổng.
  \item \textbf{Đáp án D (Nhầm lẫn số thứ tự và đơn vị đo):} Học sinh đã nhầm lẫn số thứ tự của các chữ số trong số 6 124 840 và đơn vị đo của số 800, dẫn đến việc thêm 80 và 40 vào cuối tổng.
\end{itemize}
\end{minipage}
\vspace{1em}

\begin{minipage}{\textwidth}
\noindent\textbf{Câu hỏi 8 (MEDIUM):} Số nào có chữ số 4 thuộc lớp nghìn? 
\begin{itemize}[label={}]
    \item A. 34 900
    \item B. 4 123 000
    \item C. 2 001 400
    \item D. 23 040
    \item A. Số 34 900 có chữ số 4 thuộc lớp nghìn.
    \item B. Số 4 123 000 có chữ số 4 thuộc lớp triệu.
    \item C. Số 2 001 400 có chữ số 4 thuộc lớp đơn vị.
    \item D. Số 23 040 có chữ số 4 thuộc lớp đơn vị.
\end{itemize}
\textbf{Đáp án đúng:} A \\
\textbf{Lời giải chi tiết:} A.Số 34 900 có chữ số 4 thuộc lớp nghìn. B.Số 4 123 000 có chữ số 4 thuộc lớp triệu. C.Số 2 001 400 có chữ số 4 thuộc lớp đơn vị. D.Số 23 040 có chữ số 4 thuộc lớp đơn vị. \\
\textbf{Lỗi thường gặp:}
\begin{itemize}
  \item \textbf{Đáp án B (Nhầm lẫn vị trí chữ số 4):} Học sinh nhầm lẫn vị trí chữ số 4 trong số 4 123 000, nghĩ rằng nó thuộc lớp đơn vị chứ không phải lớp triệu. Để tránh nhầm lẫn này, học sinh cần chú ý vị trí của mỗi chữ số trong số và xác định đúng lớp của mỗi chữ số.
  \item \textbf{Đáp án C (Nhầm lẫn vị trí chữ số 4):} Học sinh lại nhầm lẫn vị trí chữ số 4 trong số 2 001 400, nghĩ rằng nó thuộc lớp đơn vị chứ không phải lớp trăm. Để tránh nhầm lẫn này, học sinh cần chú ý vị trí của mỗi chữ số trong số và xác định đúng lớp của mỗi chữ số.
  \item \textbf{Đáp án D (Nhầm lẫn vị trí chữ số 4):} Học sinh lại nhầm lẫn vị trí chữ số 4 trong số 23 040, nghĩ rằng nó thuộc lớp đơn vị chứ không phải lớp nghìn. Để tránh nhầm lẫn này, học sinh cần chú ý vị trí của mỗi chữ số trong số và xác định đúng lớp của mỗi chữ số.
\end{itemize}
\end{minipage}
\vspace{1em}

\begin{minipage}{\textwidth}
\noindent\textbf{Câu hỏi 9 (MEDIUM):} Chữ số 7 trong số 123 700 405 thuộc lớp nào? 
\begin{itemize}[label={}]
    \item A. Triệu.
    \item B. Nghìn.
    \item C. Trăm.
    \item D. Đơn vị.
\end{itemize}
\textbf{Đáp án đúng:} B \\
\textbf{Lời giải chi tiết:} Chữ số 7 trong số 123 700 405 thuộc hàng trăm nghìn,lớp nghìn. \\
\textbf{Lỗi thường gặp:}
\begin{itemize}
  \item \textbf{Đáp án A (Nhầm lẫn hàng số):} Học sinh nhầm lẫn chữ số 7 thuộc hàng triệu, vì họ không phân biệt rõ ràng giữa hàng triệu và hàng trăm nghìn.
  \item \textbf{Đáp án C (Nhầm lẫn vị trí chữ số):} Học sinh nhầm lẫn vị trí chữ số 7, họ nghĩ nó thuộc hàng trăm, vì họ không hiểu rõ về cấu trúc số thập phân.
  \item \textbf{Đáp án D (Nhầm lẫn đơn vị đo):} Học sinh nhầm lẫn đơn vị đo của chữ số 7, họ nghĩ nó thuộc đơn vị, vì họ không hiểu rõ về các đơn vị đo khác nhau trong số thập phân.
\end{itemize}
\end{minipage}
\vspace{1em}

\begin{minipage}{\textwidth}
\noindent\textbf{Câu hỏi 10 (MEDIUM):} Số gồm 4 chục triệu 2 chục nghìn 7 chục 3 đơn vị viết là 
\begin{itemize}[label={}]
    \item A. 40 200 073.
    \item B. 40 020 730.
    \item C. 40 020 073
    \item D. 40 020 703
\end{itemize}
\textbf{Đáp án đúng:} C \\
\textbf{Lời giải chi tiết:} Số gồm 4 chục triệu 2 chục nghìn 7 chục 3 đơn vị được viết là:40 020 073. \\
\textbf{Lỗi thường gặp:}
\begin{itemize}
  \item \textbf{Đáp án A (Nhầm lẫn vị trí hàng chục triệu và hàng chục nghìn):} Học sinh đã nhầm lẫn vị trí hàng chục triệu và hàng chục nghìn, dẫn đến việc viết sai số. Để tránh nhầm lẫn này, học sinh cần đọc kỹ và hiểu rõ vị trí của từng hàng số trong số viết bằng chữ.
  \item \textbf{Đáp án B (Nhầm lẫn vị trí hàng chục nghìn và hàng chục):} Học sinh đã nhầm lẫn vị trí hàng chục nghìn và hàng chục, dẫn đến việc viết sai số. Để tránh nhầm lẫn này, học sinh cần đọc kỹ và hiểu rõ vị trí của từng hàng số trong số viết bằng chữ.
  \item \textbf{Đáp án D (Nhầm lẫn vị trí hàng đơn vị và hàng chục):} Học sinh đã nhầm lẫn vị trí hàng đơn vị và hàng chục, dẫn đến việc viết sai số. Để tránh nhầm lẫn này, học sinh cần đọc kỹ và hiểu rõ vị trí của từng hàng số trong số viết bằng chữ.
\end{itemize}
\end{minipage}
\vspace{1em}

\subsection{Chủ đề 1. Số tự nhiên: 13. Bài 13: Viết số tự nhiên trong hệ thập phân}
\begin{minipage}{\textwidth}
\noindent\textbf{Câu hỏi 1 (MEDIUM):} Số 67 821 đọc là: 
\begin{itemize}[label={}]
    \item A. Sáu bảy nghìn tám trăm hai mốt.
    \item B. Sáu bảy nghìn tám trăm hai mươi mốt.
    \item C. Sáu mươi bảy nghìn tám trăm hai mươi mốt.
    \item D. Sáu bảy tám trăm hai mốt.
\end{itemize}
\textbf{Đáp án đúng:} C \\
\textbf{Lời giải chi tiết:} Số 67 821 đọc là:Sáu mươi bảy nghìn tám trăm hai mươi mốt \\
\textbf{Lỗi thường gặp:}
\begin{itemize}
  \item \textbf{Đáp án A (Nhầm lẫn vị trí số):} Học sinh đã nhầm lẫn vị trí số 6 và 7, đọc số 67 821 như Sáu bảy nghìn tám trăm hai mươi mốt. Để tránh nhầm lẫn này, học sinh cần đọc số từ trái sang phải và phân biệt rõ giữa số 6 và 7.
  \item \textbf{Đáp án B (Nhầm lẫn vị trí số):} Học sinh đã nhầm lẫn vị trí số 6 và 7, đọc số 67 821 như Sáu bảy nghìn tám trăm hai mươi mốt. Để tránh nhầm lẫn này, học sinh cần đọc số từ trái sang phải và phân biệt rõ giữa số 6 và 7.
  \item \textbf{Đáp án D (Nhầm lẫn số và chữ):} Học sinh đã nhầm lẫn số 67 821 với số 6 782 1, đọc số 67 821 như Sáu bảy tám trăm hai mươi mốt. Để tránh nhầm lẫn này, học sinh cần đọc số từ trái sang phải và phân biệt rõ giữa số và chữ.
\end{itemize}
\end{minipage}
\vspace{1em}

\begin{minipage}{\textwidth}
\noindent\textbf{Câu hỏi 2 (MEDIUM):} Số “Hai trăm ba mươi bảy nghìn bốn trăm” được viết là: 
\begin{itemize}[label={}]
    \item A. 237 400 000
    \item B. 237 400
    \item C. 237
    \item D. 23 740
\end{itemize}
\textbf{Đáp án đúng:} B \\
\textbf{Lời giải chi tiết:} Số “Hai trăm ba mươi bảy nghìn bốn trăm” viết là:237 400. \\
\textbf{Lỗi thường gặp:}
\begin{itemize}
  \item \textbf{Đáp án A (Nhầm lẫn vị trí hàng nghìn):} Học sinh có thể nhầm lẫn vị trí hàng nghìn và viết sai số. Họ có thể nghĩ rằng 'Hai trăm ba mươi bảy nghìn bốn trăm' có nghĩa là 237 400 000, nhưng thực tế là hàng nghìn nên được viết sau số 400.
  \item \textbf{Đáp án C (Nhầm đơn vị đo):} Học sinh có thể nhầm lẫn đơn vị đo và viết sai số. Họ có thể nghĩ rằng 'Hai trăm ba mươi bảy nghìn bốn trăm' chỉ là 237, nhưng thực tế là số này bao gồm cả hàng nghìn và hàng trăm.
  \item \textbf{Đáp án D (Nhầm lẫn vị trí hàng trăm):} Học sinh có thể nhầm lẫn vị trí hàng trăm và viết sai số. Họ có thể nghĩ rằng 'Hai trăm ba mươi bảy nghìn bốn trăm' có nghĩa là 23 740, nhưng thực tế là hàng trăm nên được viết trước số 400.
\end{itemize}
\end{minipage}
\vspace{1em}

\begin{minipage}{\textwidth}
\noindent\textbf{Câu hỏi 3 (MEDIUM):} Viết số 28 467 000 thành tổng. 
\begin{itemize}[label={}]
    \item A. 20 000 000 + 8 000 000 + 400 000 + 60 000 + 7 000
    \item B. 20 000 000 + 8 000 000 + 400 000 + 6 000 + 700
    \item C. 20 000 000 + 8 000 000 + 400 000 + 6 000 + 7 000
    \item D. 20 000 000 + 8 000 000 + 400+ 600 + 700
\end{itemize}
\textbf{Đáp án đúng:} A \\
\textbf{Lời giải chi tiết:} Ta có: 28 467 =20 000 000 + 8 000 000 + 400 000 + 60 000 + 7 000. \\
\textbf{Lỗi thường gặp:}
\begin{itemize}
  \item \textbf{Đáp án B (Nhầm lẫn vị trí chữ số):} Học sinh có thể nhầm lẫn vị trí của chữ số 6 và 7, dẫn đến việc viết sai số 60 000 thành 6 000.
  \item \textbf{Đáp án C (Nhầm lẫn vị trí chữ số):} Học sinh có thể nhầm lẫn vị trí của chữ số 6 và 7, dẫn đến việc viết sai số 60 000 thành 7 000.
  \item \textbf{Đáp án D (Nhầm lẫn đơn vị đo):} Học sinh có thể nhầm lẫn đơn vị đo của số 400 000, dẫn đến việc viết sai số này thành 400.
\end{itemize}
\end{minipage}
\vspace{1em}

\begin{minipage}{\textwidth}
\noindent\textbf{Câu hỏi 4 (MEDIUM):} Tổng 200 000 000 + 3 000 000 + 200 000 + 50 + 1 được viết thành: 
\begin{itemize}[label={}]
    \item A. 203 200 051
    \item B. 200 320 051
    \item C. 200 320 510
    \item D. 203 002 510
\end{itemize}
\textbf{Đáp án đúng:} A \\
\textbf{Lời giải chi tiết:} Ta có: 200 000 000 + 3 000 000 + 200 000 + 50 + 1 =203 200 051 \\
\textbf{Lỗi thường gặp:}
\begin{itemize}
  \item \textbf{Đáp án B (Nhầm lẫn vị trí con số):} Học sinh đã nhầm lẫn vị trí của số 3 000 000 và số 200 000 000, dẫn đến việc cộng sai. Để tránh nhầm lẫn này, học sinh cần đọc số một cách cẩn thận và thực hiện phép tính theo thứ tự từ trái sang phải.
  \item \textbf{Đáp án C (Nhầm lẫn đơn vị đo):} Học sinh đã nhầm lẫn đơn vị đo của số 200 000 000 và số 3 000 000, dẫn đến việc cộng sai. Để tránh nhầm lẫn này, học sinh cần đọc số một cách cẩn thận và đảm bảo rằng đơn vị đo của mỗi số được sử dụng đúng.
  \item \textbf{Đáp án D (Nhầm lẫn vị trí con số và đơn vị đo):} Học sinh đã nhầm lẫn vị trí của số 3 000 000 và số 200 000 000, đồng thời cũng nhầm lẫn đơn vị đo của số 200 000 000 và số 3 000 000. Để tránh nhầm lẫn này, học sinh cần đọc số một cách cẩn thận và thực hiện phép tính theo thứ tự từ trái sang phải, đồng thời đảm bảo rằng đơn vị đo của mỗi số được sử dụng đúng.
\end{itemize}
\end{minipage}
\vspace{1em}

\begin{minipage}{\textwidth}
\noindent\textbf{Câu hỏi 5 (MEDIUM):} Số liền trước của số 87 300 là: 
\begin{itemize}[label={}]
    \item A. 87 301
    \item B. 87 289
    \item C. 87 299
    \item D. 87 310
\end{itemize}
\textbf{Đáp án đúng:} C \\
\textbf{Lời giải chi tiết:} - Muốn tìm số liền trước của một số, ta lấy số đó trừ đi 1 đơn vị. Vậy số liền trước của số 87 300 là: 87 300 – 1 =87 299 \\
\textbf{Lỗi thường gặp:}
\begin{itemize}
  \item \textbf{Đáp án A (Nhầm lẫn số thứ tự):} Học sinh có thể nhầm lẫn số thứ tự của số 87 300 và số 87 301, dẫn đến việc chọn đáp án A. Họ có thể nghĩ rằng số liền trước của 87 300 là 87 301 vì họ đã thêm 1 vào số 87 300.
  \item \textbf{Đáp án D (Nhầm lẫn số thứ tự và phép tính):} Học sinh có thể nhầm lẫn số thứ tự của số 87 300 và số 87 310, đồng thời cũng nhầm lẫn phép tính trừ đi 1 đơn vị. Họ có thể nghĩ rằng số liền trước của 87 300 là 87 310 vì họ đã thêm 10 vào số 87 300 và sau đó trừ đi 1 đơn vị.
\end{itemize}
\end{minipage}
\vspace{1em}

\begin{minipage}{\textwidth}
\noindent\textbf{Câu hỏi 6 (MEDIUM):} 10 đơn vị hay còn gọi là: 
\begin{itemize}[label={}]
    \item A. 1 chục.
    \item B. 1 trăm.
    \item C. 10 trăm
    \item D. 10 chục.
\end{itemize}
\textbf{Đáp án đúng:} A \\
\textbf{Lời giải chi tiết:} 10 đơn vị hay còn gọi là:1 chục \\
\textbf{Lỗi thường gặp:}
\begin{itemize}
  \item \textbf{Đáp án B (Nhầm lẫn đơn vị đo):} Học sinh nhầm lẫn giữa đơn vị đo 'chục' và 'trăm'. Họ nghĩ rằng 10 đơn vị là 1 trăm, nhưng thực tế 1 trăm là 10 chục.
  \item \textbf{Đáp án C (Nhầm lẫn đơn vị đo):} Học sinh nhầm lẫn giữa đơn vị đo 'trăm' và 'chục'. Họ nghĩ rằng 10 đơn vị là 10 trăm, nhưng thực tế 1 trăm là 10 chục.
  \item \textbf{Đáp án D (Nhầm lẫn đơn vị đo):} Học sinh nhầm lẫn giữa đơn vị đo 'chục' và 'trăm'. Họ nghĩ rằng 10 đơn vị là 10 chục, nhưng thực tế 1 chục là 10 đơn vị.
\end{itemize}
\end{minipage}
\vspace{1em}

\begin{minipage}{\textwidth}
\noindent\textbf{Câu hỏi 7 (MEDIUM):} Giá trị của chữ số 6 trong số 126 009 347 là: 
\begin{itemize}[label={}]
    \item A. 6 000 000
    \item B. 600 000
    \item C. 60 000
    \item D. 6 000
\end{itemize}
\textbf{Đáp án đúng:} A \\
\textbf{Lời giải chi tiết:} Giá trị của chữ số 6 trong số 126 009 347 là:6 000 000. \\
\textbf{Lỗi thường gặp:}
\begin{itemize}
  \item \textbf{Đáp án B (Nhầm lẫn vị trí chữ số 6):} Học sinh nhầm lẫn vị trí chữ số 6 trong số 126 009 347, dẫn đến việc tính giá trị của chữ số 6 là 600 000. Để làm đúng, học sinh cần xác định vị trí chữ số 6 là hàng triệu và tính giá trị của nó là 6 000 000.
  \item \textbf{Đáp án C (Nhầm lẫn đơn vị đo):} Học sinh nhầm lẫn đơn vị đo của giá trị chữ số 6, dẫn đến việc tính giá trị của chữ số 6 là 60 000. Để làm đúng, học sinh cần xác định đơn vị đo của giá trị chữ số 6 là hàng chục nghìn và tính giá trị của nó là 6 000 000.
  \item \textbf{Đáp án D (Nhầm lẫn vị trí chữ số 6):} Học sinh nhầm lẫn vị trí chữ số 6 trong số 126 009 347, dẫn đến việc tính giá trị của chữ số 6 là 6 000. Để làm đúng, học sinh cần xác định vị trí chữ số 6 là hàng trăm nghìn và tính giá trị của nó là 6 000 000.
\end{itemize}
\end{minipage}
\vspace{1em}

\begin{minipage}{\textwidth}
\noindent\textbf{Câu hỏi 8 (MEDIUM):} Trong các số 34 578; 67 900; 100 457; 9 870. Số lớn nhất là: 
\begin{itemize}[label={}]
    \item A. 34 578
    \item B. 67 900
    \item C. 100 457
    \item D. 9 870
\end{itemize}
\textbf{Đáp án đúng:} C \\
\textbf{Lời giải chi tiết:} So sánh: 9 870 < 34 578 < 67 900 < 100 457 Vậy số lớn nhất là: 100 457. \\
\textbf{Lỗi thường gặp:}
\begin{itemize}
  \item \textbf{Đáp án A (Nhầm lẫn thứ tự so sánh):} Học sinh có thể nhầm lẫn thứ tự so sánh các số, ví dụ: 34 578 < 67 900 < 100 457. Học sinh cần lưu ý so sánh từng cặp số một, thay vì so sánh tất cả các số cùng một lúc.
  \item \textbf{Đáp án B (Nhầm lẫn so sánh):} Học sinh có thể nhầm lẫn so sánh giữa 67 900 và 100 457, ví dụ: họ nghĩ 67 900 lớn hơn 100 457. Học sinh cần lưu ý so sánh từng cặp số một và sử dụng các công cụ so sánh như biểu đồ hoặc bảng so sánh.
  \item \textbf{Đáp án D (Nhầm lẫn với số nhỏ):} Học sinh có thể nhầm lẫn số 9 870 với số lớn nhất, ví dụ: họ nghĩ 9 870 là số lớn nhất. Học sinh cần lưu ý so sánh từng cặp số một và loại bỏ các số nhỏ hơn.
\end{itemize}
\end{minipage}
\vspace{1em}

\begin{minipage}{\textwidth}
\noindent\textbf{Câu hỏi 9 (MEDIUM):} Dãy số 37; 39; .....; 43; 45. Số cần điền là: 
\begin{itemize}[label={}]
    \item A. 40
    \item B. 41
    \item C. 42
    \item D. 44
\end{itemize}
\textbf{Đáp án đúng:} B \\
\textbf{Lời giải chi tiết:} - Dãy số trên gồm các số liên tiếp cách nhau 2 đơn vị. Số cần điền là: 39 + 2 =41. \\
\textbf{Lỗi thường gặp:}
\begin{itemize}
  \item \textbf{Đáp án A (Nhầm lẫn số tăng dần):} Học sinh có thể nhầm lẫn số tăng dần và số chẵn lẻ. Họ có thể nghĩ rằng số cần điền là số chẵn tiếp theo sau 39, vì vậy họ chọn đáp án 40.
  \item \textbf{Đáp án D (Nhầm lẫn số tăng dần):} Học sinh có thể nhầm lẫn số tăng dần và số chẵn lẻ. Họ có thể nghĩ rằng số cần điền là số chẵn tiếp theo sau 43, vì vậy họ chọn đáp án 44.
\end{itemize}
\end{minipage}
\vspace{1em}

\begin{minipage}{\textwidth}
\noindent\textbf{Câu hỏi 10 (MEDIUM):} Điền dấu thích hợp vào chỗ chấm: \\
3 478 ..... 3 000 + 400 + 70 + 8 
\begin{itemize}[label={}]
    \item A. <
    \item B. >
    \item C. =
\end{itemize}
\textbf{Đáp án đúng:} C \\
\textbf{Lời giải chi tiết:} Ta có: 3 000 + 400 + 70 + 8 = 3 478 So sánh: 3 478 = 3 478 hay 3 478=3 000 + 400 + 70 + 8 \\
\textbf{Lỗi thường gặp:}
\begin{itemize}
  \item \textbf{Đáp án A (Nhầm lẫn thứ tự cộng trừ):} Học sinh có thể nhầm lẫn thứ tự cộng trừ, ví dụ: họ nghĩ rằng 3 478 lớn hơn 3 000 nên chọn đáp án A. Để tránh nhầm lẫn này, học sinh cần đọc kỹ câu hỏi và thực hiện phép tính một cách cẩn thận.
  \item \textbf{Đáp án B (Nhầm lẫn về phép tính):} Học sinh có thể nhầm lẫn về phép tính, ví dụ: họ nghĩ rằng 3 478 lớn hơn 3 000 + 400 + 70 + 8 nên chọn đáp án B. Để tránh nhầm lẫn này, học sinh cần thực hiện phép tính một cách cẩn thận và kiểm tra lại kết quả.
\end{itemize}
\end{minipage}
\vspace{1em}

\subsection{Chủ đề 2. Các phép tính với số tự nhiên: 17. Bài 42: Chia cho số có hai chữ số (tiếp theo)}
\begin{minipage}{\textwidth}
\noindent\textbf{Câu hỏi 1 (MEDIUM):} Làm tròn đến hàng chục số bị chia. 
\begin{itemize}[label={}]
    \item A. 230
    \item B. 240
    \item C. 200
    \item D. 300
\end{itemize}
\textbf{Đáp án đúng:} B \\
\textbf{Lời giải chi tiết:} Số bị chia là: 236 Số 236 có chữ số hàng đơn vị là 6 > 5, nên ta làm tròn lên Vậy: Làm tròn số 236 đến hàng chục được240 \\
\textbf{Lỗi thường gặp:}
\begin{itemize}
  \item \textbf{Đáp án A (Nhầm lẫn chữ số hàng đơn vị):} Học sinh nhầm lẫn chữ số hàng đơn vị là 3 với chữ số hàng đơn vị là 6, dẫn đến việc chọn đáp án 230.
  \item \textbf{Đáp án C (Nhầm lẫn số bị chia):} Học sinh nhầm lẫn số bị chia là 236 với số khác, dẫn đến việc chọn đáp án 200.
  \item \textbf{Đáp án D (Nhầm lẫn quy tắc làm tròn):} Học sinh nhầm lẫn quy tắc làm tròn số bị chia, dẫn đến việc chọn đáp án 300.
\end{itemize}
\end{minipage}
\vspace{1em}

\begin{minipage}{\textwidth}
\noindent\textbf{Câu hỏi 2 (MEDIUM):} Làm tròn đến hàng chục số chia. 
\begin{itemize}[label={}]
    \item A. 60
    \item B. 50
    \item C. 70
    \item D. 40
\end{itemize}
\textbf{Đáp án đúng:} A \\
\textbf{Lời giải chi tiết:} Số chia là: 58 Số 58 có chữ số hàng đơn vị là 8 > 5, nên ta làm tròn lên Vậy: Làm tròn số 58 đến hàng chục được60 \\
\textbf{Lỗi thường gặp:}
\begin{itemize}
  \item \textbf{Đáp án B (Nhầm lẫn chữ số hàng đơn vị):} Học sinh nhầm lẫn chữ số hàng đơn vị là 8 > 5 nên làm tròn lên, nhưng thực tế chỉ cần xem chữ số hàng đơn vị là 8 > 5, nên làm tròn lên hàng chục, không cần xem chữ số hàng đơn vị. Cách làm đúng: Xem chữ số hàng đơn vị là 8 > 5, nên làm tròn lên hàng chục.
  \item \textbf{Đáp án C (Nhầm lẫn về số chia):} Học sinh nhầm lẫn số chia là 58, nên làm tròn lên hàng chục, nhưng thực tế số chia là 58, không cần làm tròn lên. Cách làm đúng: Xem số chia là 58, không cần làm tròn lên.
  \item \textbf{Đáp án D (Nhầm lẫn về quy tắc làm tròn):} Học sinh nhầm lẫn quy tắc làm tròn, nên làm tròn xuống hàng chục, nhưng thực tế số chia là 58, nên làm tròn lên hàng chục. Cách làm đúng: Xem số chia là 58, nên làm tròn lên hàng chục.
\end{itemize}
\end{minipage}
\vspace{1em}

\begin{minipage}{\textwidth}
\noindent\textbf{Câu hỏi 3 (MEDIUM):} Thương ước lượng sau khi làm tròn số bị chia và số chia đến hàng chục. 
\begin{itemize}[label={}]
    \item A. 1
    \item B. 3
    \item C. 2
    \item D. 4
\end{itemize}
\textbf{Đáp án đúng:} D \\
\textbf{Lời giải chi tiết:} Ta có: 240 : 60 = 4 Vậy thương ước lượng sau khi làm tròn số bị chia và số chia đến hàng chục là: 4 \\
\textbf{Lỗi thường gặp:}
\begin{itemize}
  \item \textbf{Đáp án A (Nhầm lẫn về việc làm tròn số bị chia và số chia):} Học sinh có thể nhầm lẫn về việc làm tròn số bị chia và số chia, dẫn đến việc tính toán không chính xác. Ví dụ, họ có thể nghĩ rằng số bị chia là 240, nhưng thực tế là 240 đã được làm tròn đến hàng chục, nên số bị chia thực tế là 240 trở xuống. Điều này dẫn đến kết quả không chính xác là 1.
  \item \textbf{Đáp án B (Quên làm tròn số chia):} Học sinh có thể quên làm tròn số chia, dẫn đến việc tính toán không chính xác. Ví dụ, họ có thể nghĩ rằng số chia là 60, nhưng thực tế là 60 đã được làm tròn đến hàng chục, nên số chia thực tế là 60 trở xuống. Điều này dẫn đến kết quả không chính xác là 3.
  \item \textbf{Đáp án C (Làm tròn số bị chia không đúng):} Học sinh có thể làm tròn số bị chia không đúng, dẫn đến việc tính toán không chính xác. Ví dụ, họ có thể nghĩ rằng số bị chia là 240, nhưng thực tế là 240 đã được làm tròn đến hàng chục, nên số bị chia thực tế là 240 trở xuống. Điều này dẫn đến kết quả không chính xác là 2.
\end{itemize}
\end{minipage}
\vspace{1em}

\begin{minipage}{\textwidth}
\noindent\textbf{Câu hỏi 4 (MEDIUM):} Thương của phép chia 236 : 58 là bao nhiêu? 
\begin{itemize}[label={}]
    \item A. 2
    \item B. 3
    \item C. 4
    \item D. 5
\end{itemize}
\textbf{Đáp án đúng:} C \\
\textbf{Lời giải chi tiết:} Vậy thương của phép chia 236 : 58 là 4. \\
\textbf{Lỗi thường gặp:}
\begin{itemize}
  \item \textbf{Đáp án A (Nhầm lẫn phép chia):} Học sinh có thể nhầm lẫn phép chia 236 : 58 với phép chia 236 : 5 hoặc 236 : 2, dẫn đến việc chọn đáp án A. Để tránh nhầm lẫn này, học sinh cần đọc kỹ số chia và số bị chia trước khi thực hiện phép chia.
  \item \textbf{Đáp án B (Quên số chia):} Học sinh có thể quên số chia 58 và chỉ tập trung vào số bị chia 236, dẫn đến việc chọn đáp án B. Để tránh quên số chia, học sinh cần đọc kỹ số chia trước khi thực hiện phép chia.
  \item \textbf{Đáp án D (Tính sai):} Học sinh có thể tính sai thương của phép chia 236 : 58, dẫn đến việc chọn đáp án D. Để tránh tính sai, học sinh cần kiểm tra lại phép chia và đảm bảo rằng thương được tính chính xác.
\end{itemize}
\end{minipage}
\vspace{1em}

\begin{minipage}{\textwidth}
\noindent\textbf{Câu hỏi 5 (MEDIUM):} So sánh thương của hai phép chia sau \\
176 : 44 ..... 212 : 53 
\begin{itemize}[label={}]
    \item A. >
    \item B. <
    \item C. =
\end{itemize}
\textbf{Đáp án đúng:} C \\
\textbf{Lời giải chi tiết:} Ta có: 176 : 44= 4 212 : 53 = 4 So sánh: 4 = 4 hay 176 : 44=212 : 53 \\
\textbf{Lỗi thường gặp:}
\begin{itemize}
  \item \textbf{Đáp án A (Nhầm lẫn với phép chia không chính xác):} Học sinh có thể nhầm lẫn với phép chia không chính xác, ví dụ như 176 : 44 = 4 và 212 : 53 = 4, nhưng thực tế 212 không chia hết cho 53, nên phép chia không chính xác.
  \item \textbf{Đáp án B (Nhầm lẫn với phép chia không chính xác và không hiểu rõ khái niệm thương):} Học sinh có thể nhầm lẫn với phép chia không chính xác và không hiểu rõ khái niệm thương, ví dụ như 176 : 44 = 4 và 212 : 53 = 4, nhưng thực tế 212 không chia hết cho 53, nên phép chia không chính xác và không hiểu rõ khái niệm thương.
\end{itemize}
\end{minipage}
\vspace{1em}

\begin{minipage}{\textwidth}
\noindent\textbf{Câu hỏi 6 (MEDIUM):} Phép chia nào dưới đây có số dư bằng 4? 
\begin{itemize}[label={}]
    \item A. 240 : 48
    \item B. 310 : 34
    \item C. 175 : 25
    \item D. 132 : 33
\end{itemize}
\textbf{Đáp án đúng:} B \\
\textbf{Lời giải chi tiết:} Ta có: Vậy phép chia có số dư bằng 4 là310 : 34. \\
\textbf{Lỗi thường gặp:}
\begin{itemize}
  \item \textbf{Đáp án A (Nhầm lẫn phép chia):} Học sinh có thể nhầm lẫn phép chia 240 : 48 vì số chia lớn hơn số bị chia, dẫn đến kết quả không phải là số dư bằng 4. Để tránh nhầm lẫn này, học sinh cần kiểm tra lại phép chia và đảm bảo rằng số chia nhỏ hơn hoặc bằng số bị chia.
  \item \textbf{Đáp án C (Nhầm lẫn phép chia):} Học sinh có thể nhầm lẫn phép chia 175 : 25 vì số chia lớn hơn số bị chia, dẫn đến kết quả không phải là số dư bằng 4. Để tránh nhầm lẫn này, học sinh cần kiểm tra lại phép chia và đảm bảo rằng số chia nhỏ hơn hoặc bằng số bị chia.
  \item \textbf{Đáp án D (Nhầm lẫn phép chia):} Học sinh có thể nhầm lẫn phép chia 132 : 33 vì số chia lớn hơn số bị chia, dẫn đến kết quả không phải là số dư bằng 4. Để tránh nhầm lẫn này, học sinh cần kiểm tra lại phép chia và đảm bảo rằng số chia nhỏ hơn hoặc bằng số bị chia.
\end{itemize}
\end{minipage}
\vspace{1em}

\begin{minipage}{\textwidth}
\noindent\textbf{Câu hỏi 7 (MEDIUM):} Số dư của phép chia 327 : 14 là bao nhiêu? 
\begin{itemize}[label={}]
    \item A. 1
    \item B. 3
    \item C. 5
    \item D. 7
\end{itemize}
\textbf{Đáp án đúng:} C \\
\textbf{Lời giải chi tiết:} Vậy số dư của phép chia 327 : 14 là:5. \\
\textbf{Lỗi thường gặp:}
\begin{itemize}
  \item \textbf{Đáp án A (Nhầm lẫn phép chia không hết):} Học sinh có thể nhầm lẫn phép chia không hết với phép chia hết, dẫn đến việc chọn đáp án A. Để tránh nhầm lẫn này, học sinh cần hiểu rõ rằng phép chia không hết sẽ luôn có số dư.
  \item \textbf{Đáp án B (Nhầm lẫn số dư):} Học sinh có thể nhầm lẫn số dư là 3, vì 14 x 23 = 322 và 327 - 322 = 5. Để tránh nhầm lẫn này, học sinh cần kiểm tra lại phép chia và tính số dư chính xác.
  \item \textbf{Đáp án D (Nhầm lẫn phép chia):} Học sinh có thể nhầm lẫn phép chia 327 : 14 với phép chia 14 : 327, dẫn đến việc chọn đáp án D. Để tránh nhầm lẫn này, học sinh cần hiểu rõ rằng phép chia là một quá trình chia số lớn hơn cho số nhỏ hơn.
\end{itemize}
\end{minipage}
\vspace{1em}

\begin{minipage}{\textwidth}
\noindent\textbf{Câu hỏi 8 (MEDIUM):} Thương của phép chia 281 : 41 là: 
\begin{itemize}[label={}]
    \item A. 6
    \item B. 35
    \item C. 8
    \item D. 5
\end{itemize}
\textbf{Đáp án đúng:} A \\
\textbf{Lời giải chi tiết:} Vậy thương của 281 : 41 là: 6. \\
\textbf{Lỗi thường gặp:}
\begin{itemize}
  \item \textbf{Đáp án B (Lỗi tính toán):} Học sinh chọn đáp án B do nhầm lẫn trong các bước cộng, trừ, nhân, chia hoặc đọc sai đơn vị đo.
  \item \textbf{Đáp án C (Lỗi tính toán):} Học sinh chọn đáp án C do nhầm lẫn trong các bước cộng, trừ, nhân, chia hoặc đọc sai đơn vị đo.
  \item \textbf{Đáp án D (Lỗi tính toán):} Học sinh chọn đáp án D do nhầm lẫn trong các bước cộng, trừ, nhân, chia hoặc đọc sai đơn vị đo.
\end{itemize}
\end{minipage}
\vspace{1em}

\begin{minipage}{\textwidth}
\noindent\textbf{Câu hỏi 9 (MEDIUM):} Thương chính xác của phép chia 246 : 41 là: 
\begin{itemize}[label={}]
    \item A. 5
    \item B. 6
    \item C. 7
    \item D. 5
\end{itemize}
\textbf{Đáp án đúng:} B \\
\textbf{Lời giải chi tiết:} Vậy thương chính xác của phép chia 246 : 41 là:6. \\
\textbf{Lỗi thường gặp:}
\begin{itemize}
  \item \textbf{Đáp án A (Nhầm lẫn phép chia không chính xác):} Học sinh có thể nhầm lẫn phép chia không chính xác vì số chia lớn hơn số bị chia, dẫn đến kết quả không chính xác. Để tránh nhầm lẫn này, học sinh cần thực hiện phép chia một cách cẩn thận và kiểm tra lại kết quả.
  \item \textbf{Đáp án C (Quên số dư):} Học sinh có thể quên số dư khi thực hiện phép chia, dẫn đến kết quả không chính xác. Để tránh nhầm lẫn này, học sinh cần tính toán số dư và kiểm tra lại kết quả.
  \item \textbf{Đáp án D (Nhầm lẫn với đáp án khác):} Học sinh có thể nhầm lẫn với đáp án khác vì số chia và số bị chia có nhiều điểm tương đồng. Để tránh nhầm lẫn này, học sinh cần đọc kỹ câu hỏi và kiểm tra lại đáp án.
\end{itemize}
\end{minipage}
\vspace{1em}

\begin{minipage}{\textwidth}
\noindent\textbf{Câu hỏi 10 (MEDIUM):} Kết quả của phép tính 896 : 32 là 
\begin{itemize}[label={}]
    \item A. 64
    \item B. 28
    \item C. 56
    \item D. 25
\end{itemize}
\textbf{Đáp án đúng:} B \\
\textbf{Lời giải chi tiết:} Vậy kết quả của phép tính 896 : 32 là:28. \\
\textbf{Lỗi thường gặp:}
\begin{itemize}
  \item \textbf{Đáp án A (Nhầm lẫn về phép chia):} Học sinh có thể nhầm lẫn về phép chia, ví dụ họ nghĩ rằng 896 chia cho 32 là 28, nhưng thực tế là 896 chia cho 32 là 28, nhưng họ lại nghĩ rằng 896 chia cho 32 là 64 vì họ đã nhầm lẫn về phép chia.
  \item \textbf{Đáp án C (Nhầm lẫn về phép chia):} Học sinh có thể nhầm lẫn về phép chia, ví dụ họ nghĩ rằng 896 chia cho 32 là 28, nhưng thực tế là 896 chia cho 32 là 28, nhưng họ lại nghĩ rằng 896 chia cho 32 là 56 vì họ đã nhầm lẫn về phép chia.
  \item \textbf{Đáp án D (Nhầm lẫn về phép chia):} Học sinh có thể nhầm lẫn về phép chia, ví dụ họ nghĩ rằng 896 chia cho 32 là 28, nhưng thực tế là 896 chia cho 32 là 28, nhưng họ lại nghĩ rằng 896 chia cho 32 là 25 vì họ đã nhầm lẫn về phép chia.
\end{itemize}
\end{minipage}
\vspace{1em}

\subsection{Chủ đề 2. Các phép tính với số tự nhiên: 19. Bài 44: Thương có chữ số 0}
\begin{minipage}{\textwidth}
\noindent\textbf{Câu hỏi 1 (MEDIUM):} Trong các phép chia sau, phép chia nào mà thương có chứa chữ số 0? 
\begin{itemize}[label={}]
    \item A. 132 : 11
    \item B. 1 236 : 12
    \item C. 1 150 : 46
    \item D. 1 092 : 52
\end{itemize}
\textbf{Đáp án đúng:} B \\
\textbf{Lời giải chi tiết:} Vậy phép chia mà thương có chứa chữ số 0 là:1 236 : 12. \\
\textbf{Lỗi thường gặp:}
\begin{itemize}
  \item \textbf{Đáp án A (Nhầm lẫn số nguyên tố):} Học sinh có thể nhầm lẫn số 11 là số nguyên tố, dẫn đến việc chọn đáp án A. Tuy nhiên, số 11 không phải là yếu tố quyết định trong phép chia này. Học sinh cần tập trung vào việc tìm kiếm số chia có thể chia hết cho cả 11 và 12, như trong đáp án B.
  \item \textbf{Đáp án C (Tính sai về số chia):} Học sinh có thể tính sai về số chia 46, dẫn đến việc chọn đáp án C. Tuy nhiên, số 46 không phải là số chia có thể chia hết cho cả 150 và 46. Học sinh cần tập trung vào việc tìm kiếm số chia có thể chia hết cho cả 150 và 46, như trong đáp án B.
  \item \textbf{Đáp án D (Tính sai về số chia):} Học sinh có thể tính sai về số chia 52, dẫn đến việc chọn đáp án D. Tuy nhiên, số 52 không phải là số chia có thể chia hết cho cả 1092 và 52. Học sinh cần tập trung vào việc tìm kiếm số chia có thể chia hết cho cả 1092 và 52, như trong đáp án B.
\end{itemize}
\end{minipage}
\vspace{1em}

\begin{minipage}{\textwidth}
\noindent\textbf{Câu hỏi 2 (MEDIUM):} Đâu là kết quả của phép tính 4 725 : 45? 
\begin{itemize}[label={}]
    \item A. 103
    \item B. 104
    \item C. 105
    \item D. 106
\end{itemize}
\textbf{Đáp án đúng:} C \\
\textbf{Lời giải chi tiết:} Vậy kết quả của phép tính 4 725 : 45 là:105. \\
\textbf{Lỗi thường gặp:}
\begin{itemize}
  \item \textbf{Đáp án A (Nhầm lẫn vị trí chữ số):} Học sinh có thể nhầm lẫn vị trí chữ số 4 và 5 khi thực hiện phép tính chia, dẫn đến kết quả là 103.
  \item \textbf{Đáp án B (Nhầm lẫn vị trí chữ số):} Học sinh có thể nhầm lẫn vị trí chữ số 4 và 5 khi thực hiện phép tính chia, dẫn đến kết quả là 104.
  \item \textbf{Đáp án D (Nhầm lẫn vị trí chữ số):} Học sinh có thể nhầm lẫn vị trí chữ số 4 và 6 khi thực hiện phép tính chia, dẫn đến kết quả là 106.
\end{itemize}
\end{minipage}
\vspace{1em}

\begin{minipage}{\textwidth}
\noindent\textbf{Câu hỏi 3 (MEDIUM):} Thương của phép chia 9 588 : 47 có chữ số 0 thuộc hàng nào? 
\begin{itemize}[label={}]
    \item A. Hàng đơn vị.
    \item B. Hàng chục.
    \item C. Hàng trăm
    \item D. Không xác định.
\end{itemize}
\textbf{Đáp án đúng:} B \\
\textbf{Lời giải chi tiết:} Phép tính trên có thương là: 204. Số 204 có chữ số 0 thuộc hàng chục. Vậy thương của phép chia 9 588 : 47 có chữ số 0 thuộc hàng chục. \\
\textbf{Lỗi thường gặp:}
\begin{itemize}
  \item \textbf{Đáp án A (Nhầm lẫn vị trí chữ số 0):} Học sinh có thể nhầm lẫn vị trí chữ số 0 trong số 204, nghĩ rằng nó thuộc hàng đơn vị. Để tránh nhầm lẫn này, học sinh cần lưu ý rằng chữ số 0 thường nằm ở hàng chục hoặc hàng trăm, chứ không phải hàng đơn vị.
  \item \textbf{Đáp án C (Nhầm lẫn vị trí chữ số 0):} Học sinh có thể nghĩ rằng chữ số 0 thuộc hàng trăm, vì số 204 có chữ số 0 ở vị trí hàng trăm. Tuy nhiên, trong trường hợp này, chữ số 0 thuộc hàng chục, không phải hàng trăm.
  \item \textbf{Đáp án D (Trốn tránh việc giải thích):} Học sinh có thể chọn đáp án D để trốn tránh việc giải thích và tìm hiểu lý do tại sao thương có chữ số 0 thuộc hàng chục. Tuy nhiên, việc giải thích và tìm hiểu lý do là một phần quan trọng của quá trình học tập và giúp học sinh hiểu rõ hơn về vấn đề.
\end{itemize}
\end{minipage}
\vspace{1em}

\begin{minipage}{\textwidth}
\noindent\textbf{Câu hỏi 4 (MEDIUM):} “Số 0 ở thương của phép chia 1 350 : 27 thuộc hàng chục”. Em hãy cho biết, phát biểu trên đúng hay sai? 
\begin{itemize}[label={}]
    \item A. Đúng.
    \item B. Sai.
\end{itemize}
\textbf{Đáp án đúng:} B \\
\textbf{Lời giải chi tiết:} Phép tính trên có thương là 50. Số 50 có chữ số 0 thuộc hàng đơn vị. Vậy phát biểu trên là sai. \\
\textbf{Lỗi thường gặp:}
\begin{itemize}
  \item \textbf{Đáp án A (Nhầm lẫn vị trí chữ số 0):} Học sinh có thể nhầm lẫn vị trí chữ số 0 trong số 0 ở hàng chục và số 0 ở hàng đơn vị. Để tránh nhầm lẫn này, học sinh cần chú ý đến vị trí của chữ số 0 trong mỗi hàng.
  \item \textbf{Đáp án B (Nhầm lẫn vị trí chữ số 0):} Học sinh có thể nhầm lẫn vị trí chữ số 0 trong số 0 ở hàng chục và số 0 ở hàng đơn vị. Để tránh nhầm lẫn này, học sinh cần chú ý đến vị trí của chữ số 0 trong mỗi hàng.
\end{itemize}
\end{minipage}
\vspace{1em}

\begin{minipage}{\textwidth}
\noindent\textbf{Câu hỏi 5 (MEDIUM):} Thương của phép tính 5 250 : 35 là? 
\begin{itemize}[label={}]
    \item A. 15
    \item B. 150
    \item C. 105
    \item D. 140
\end{itemize}
\textbf{Đáp án đúng:} B \\
\textbf{Lời giải chi tiết:} Vậy thương của phép tính 5 250 : 35 là 150. \\
\textbf{Lỗi thường gặp:}
\begin{itemize}
  \item \textbf{Đáp án A (Nhầm lẫn đơn vị đo):} Học sinh có thể nhầm lẫn đơn vị đo của số chia (35) và số bị chia (250), dẫn đến việc tính ra đáp án 15. Để tránh nhầm lẫn này, học sinh cần đọc kỹ và hiểu rõ đơn vị đo của từng số.
  \item \textbf{Đáp án C (Nhầm lẫn phép chia):} Học sinh có thể nhầm lẫn phép chia 5250 : 35 với phép chia 525 : 3,5, dẫn đến việc tính ra đáp án 105. Để tránh nhầm lẫn này, học sinh cần đọc kỹ và hiểu rõ phép chia cần thực hiện.
  \item \textbf{Đáp án D (Nhầm lẫn số chia):} Học sinh có thể nhầm lẫn số chia (35) với số bị chia (250), dẫn đến việc tính ra đáp án 140. Để tránh nhầm lẫn này, học sinh cần đọc kỹ và hiểu rõ số chia và số bị chia.
\end{itemize}
\end{minipage}
\vspace{1em}

\begin{minipage}{\textwidth}
\noindent\textbf{Câu hỏi 6 (MEDIUM):} Kết quả của phép tính 3 300 : 22 là: 
\begin{itemize}[label={}]
    \item A. 15
    \item B. 105
    \item C. 150
    \item D. 120
\end{itemize}
\textbf{Đáp án đúng:} C \\
\textbf{Lời giải chi tiết:} Vậy kết quả của 3 300 : 22 là: 150. \\
\textbf{Lỗi thường gặp:}
\begin{itemize}
  \item \textbf{Đáp án A (Nhầm lẫn phép chia với phép nhân):} Học sinh có thể nhầm lẫn phép chia với phép nhân, vì họ đã từng thực hiện phép nhân 3 và 300 để nhận được kết quả 900, sau đó nhầm lẫn và chọn đáp án A là 15.
  \item \textbf{Đáp án B (Nhầm lẫn đơn vị đo):} Học sinh có thể nhầm lẫn đơn vị đo và chọn đáp án B là 105, vì họ đã từng thực hiện phép chia 3000 với 22 để nhận được kết quả 136,36, sau đó nhầm lẫn và chọn đáp án B.
  \item \textbf{Đáp án D (Nhầm lẫn kết quả phép chia):} Học sinh có thể nhầm lẫn kết quả phép chia và chọn đáp án D là 120, vì họ đã từng thực hiện phép chia 3000 với 22 để nhận được kết quả 136,36, sau đó nhầm lẫn và chọn đáp án D.
\end{itemize}
\end{minipage}
\vspace{1em}

\begin{minipage}{\textwidth}
\noindent\textbf{Câu hỏi 7 (MEDIUM):} Kết quả của phép tính 3 780 : 35 là: 
\begin{itemize}[label={}]
    \item A. 108
    \item B. 180
    \item C. 18
    \item D. 150
\end{itemize}
\textbf{Đáp án đúng:} A \\
\textbf{Lời giải chi tiết:} Vậy số cần điền là: 108. \\
\textbf{Lỗi thường gặp:}
\begin{itemize}
  \item \textbf{Đáp án B (Nhầm lẫn vị trí chữ số):} Học sinh có thể nhầm lẫn vị trí chữ số 8 và 0, dẫn đến việc tính ra đáp án 180. Để tránh nhầm lẫn này, học sinh cần đọc số một cách cẩn thận và đảm bảo rằng mỗi chữ số được đặt đúng vị trí.
  \item \textbf{Đáp án C (Nhầm lẫn đơn vị đo):} Học sinh có thể nhầm lẫn đơn vị đo của phép tính, dẫn đến việc tính ra đáp án 18. Để tránh nhầm lẫn này, học sinh cần đọc kỹ đơn vị đo và đảm bảo rằng nó phù hợp với phép tính.
  \item \textbf{Đáp án D (Nhầm lẫn phép tính):} Học sinh có thể nhầm lẫn phép tính 780 : 35 với phép tính 780 - 35, dẫn đến việc tính ra đáp án 150. Để tránh nhầm lẫn này, học sinh cần đọc kỹ phép tính và đảm bảo rằng nó được thực hiện đúng.
\end{itemize}
\end{minipage}
\vspace{1em}

\begin{minipage}{\textwidth}
\noindent\textbf{Câu hỏi 8 (MEDIUM):} 4 905 : ..... = 45 
\begin{itemize}[label={}]
    \item A. 107
    \item B. 108
    \item C. 109
    \item D. 110
\end{itemize}
\textbf{Đáp án đúng:} C \\
\textbf{Lời giải chi tiết:} - Muốn tìm số chia, ta lấy số bị chia chia cho thương. Vậy số cần điền là 109. \\
\textbf{Lỗi thường gặp:}
\begin{itemize}
  \item \textbf{Đáp án A (Nhầm lẫn công thức chia):} Học sinh có thể nhầm lẫn công thức chia số với công thức tìm số chia. Họ có thể tính 4.905 ÷ 45 = 109, nhưng thực tế họ đang tìm số chia, không phải thương.
  \item \textbf{Đáp án B (Quên đổi dấu):} Học sinh có thể quên đổi dấu khi tính số chia. Họ có thể tính 4.905 ÷ 45 = -109, nhưng thực tế số chia luôn là số nguyên dương.
  \item \textbf{Đáp án D (Tính sai số chia):} Học sinh có thể tính sai số chia do nhầm lẫn công thức chia số với công thức tìm số chia. Họ có thể tính 4.905 ÷ 45 = 110, nhưng thực tế số chia không phải là 110.
\end{itemize}
\end{minipage}
\vspace{1em}

\begin{minipage}{\textwidth}
\noindent\textbf{Câu hỏi 9 (MEDIUM):} ..... : 56 = 1 700 : 25 
\begin{itemize}[label={}]
    \item A. 68
    \item B. 1 400
    \item C. 3 808
    \item D. 1 619
\end{itemize}
\textbf{Đáp án đúng:} C \\
\textbf{Lời giải chi tiết:} Số cần tìm là số bị chia Số chia là: 56 Thương là kết quả của phép tính 1 700 : 25 Ta có: 1 700 : 25 = 68 - Muốn tìm số bị chia, ta lấy thương nhân với số chia. Số bị chia là: 68 × 56 = 3 808 Vậy số cần điền là:3 808. \\
\textbf{Lỗi thường gặp:}
\begin{itemize}
  \item \textbf{Đáp án A (Nhầm lẫn công thức):} Học sinh có thể nhầm lẫn công thức tính thương, dẫn đến việc lấy 1 700 : 25 = 68 và nhân với số chia 56, thay vì nhân số chia 25 với thương 68.
  \item \textbf{Đáp án B (Lỗi đơn vị đo):} Học sinh có thể nhầm lẫn đơn vị đo, lấy 1 700 : 25 = 68 và nhân với số chia 56, thay vì nhân số chia 25 với thương 68.
  \item \textbf{Đáp án D (Lỗi nhân số):} Học sinh có thể nhầm lẫn việc nhân số, lấy 1 700 : 25 = 68 và nhân với số chia 56, thay vì nhân số chia 25 với thương 68.
\end{itemize}
\end{minipage}
\vspace{1em}

\begin{minipage}{\textwidth}
\noindent\textbf{Câu hỏi 10 (MEDIUM):} Điền dấu thích hợp vào chỗ chấm. \\
13 × 11 ..... 1 800 : 12 
\begin{itemize}[label={}]
    \item A. <
    \item B. >
    \item C. =
\end{itemize}
\textbf{Đáp án đúng:} A \\
\textbf{Lời giải chi tiết:} 13 × 11 = 143 1 800 : 12 = 150 So sánh: 143 < 150 hay 13 × 11<1 800 : 12 \\
\textbf{Lỗi thường gặp:}
\begin{itemize}
  \item \textbf{Đáp án C (Nhầm lẫn về dấu hiệu so sánh):} Học sinh có thể nhầm lẫn giữa dấu hiệu so sánh và dấu hiệu bằng. Họ có thể nghĩ rằng 13 × 11 = 1 800 : 12 vì cả hai biểu thức đều có kết quả bằng nhau, nhưng thực tế là 13 × 11 < 1 800 : 12.
  \item \textbf{Đáp án B (Nhầm lẫn về đơn vị đo):} Học sinh có thể nhầm lẫn giữa đơn vị đo của hai biểu thức. Họ có thể nghĩ rằng 13 × 11 = 1 800 : 12 vì cả hai biểu thức đều có kết quả bằng nhau, nhưng thực tế là 13 × 11 < 1 800 : 12.
\end{itemize}
\end{minipage}
\vspace{1em}

\end{document}
